\section{Introduction}

Curve-counting invariants of a Calabi--Yau threefold change under a
conifold transition by summing over curve classes on the small
resolution. The numerical formula holds in every genus, but its
interpretation in a universal coefficient algebra requires a single
support bound valid before any numerical evaluation. Such a bound
must also account for the whole exceptional contribution and specify
the completion in which specialization is performed.

We combine Pardon's universal curve-counting algebra with the
numerical conifold formula to obtain this interpretation, and then
compute every source contribution in one nonexceptional class of a
compact example. The geometric step in the computation is an
Euler-characteristic bound for pure curves, including nonreduced
exceptional components. It makes a finite genus-zero calculation
sufficient to determine the full universal coefficients.

For the first result, let $Y\to X_0\rightsquigarrow S$ be a
projective conifold transition with trivial canonical bundles,
torsion-free second integral homology, and a smoothing whose total
space is a smooth algebraic fourfold. Fix the common-complement
marking, giving $\pi:H_2(Y;\ZZ)\to H_2(S;\ZZ)$. Pardon's local
generators $x_{g,m}$ form the rational polynomial coefficient
algebra $H$, and $f_g(z)=1+\sum_{m\ge1}x_{g,m}z^m$ is the local
series. Write $\Z_V$ for the universal effective-cycle series of a
threefold $V$. Theorem~\ref{thm:specialization} proves
\[
 \Z_S=\Sp_\pi\left(\frac{\Z_Y}{\prod_i f_0(Q^{[C_i]})}\right),
\]
where $C_i$ are the exceptional curves. Every coefficient is a finite
sum of polynomials in $H$. The proof uses an ample divisor on the
common blow-up to bound all effective lifts in bounded target degree.
Empty-contact degeneration then controls universal logarithmic
support. We apply the total-space-class form of degeneration to
handle projected classes which might be ineffective on the smoothing.

The compact application uses a four-node toric hypersurface
transition, specified in Section~\ref{sec:geometry} and
Appendix~\ref{app:data}. Its smoothing contains a plane whose
primitive line class is $\beta=(1,1,0)$. In a fixed integral source
basis, $\pi(n_1,\ldots,n_6)=(n_2,n_3,n_4)$. There are 37 possible
normalized contributions over $\beta$ under our finite support
bound. Section~\ref{sec:euler} proves that every pure curve in the
three bidegrees needed for this calculation has Euler characteristic
at least one, even after adding arbitrary multiple exceptional
components. The proof identifies the contracted locus as two
rational surfaces and two isolated rational curves, then bounds
intersection lengths with multiple $(-1,-1)$-curves. This removes
all positive-genus primitive generators from the relevant universal
logarithm. Section~\ref{sec:calculation} consequently determines
every one of the 37 coefficients, of which eight are nonzero.
Their sum is $3x_{0,1}$.

An independent stable-pair calculation on the smoothing plane gives
\[
 U_{S,\beta}=3x_{0,1},\qquad
 U_{S,2\beta}=3x_{0,2}+3x_{0,1}^2-6x_{0,1}.
\]
These are identities in $H$. The Euler-characteristic arguments,
rather than agreement of finitely many GW invariants, exclude the
remaining primitive generators.

Pardon's structure theorem and connected evaluations are the algebraic
inputs throughout \cite{Pardon}. The all-genus numerical conifold
formula and its finite class sums are already established by
Lee--Lin--Wang \cite[Proposition 2.1]{LLW}. Lin--Wang give
disconnected GW and PT transition formulas with exceptional factors
\cite[Proposition 3.5]{LW}, and a transfer theorem for stationary
descendent correspondence. Our unmarked statement does not extend
that descendent theorem. Its content is the coefficientwise finite
identity in the universal algebra, together with the complete
source calculation just described. The values $3$ and $-6$ on the
plane are familiar local-curve-counting values; their role here is
to determine universal polynomials and to check the compact
specialization independently.

The product exponents obey a finite sum formula,
Corollary~\ref{cor:integral}, so integrality on the resolution implies
integrality on the smoothing. The hypothesis is substantive: the
scalar equivariant local generators do not give integer exponents
for every projective Calabi--Yau threefold. The companion paper
\cite{WuebbenIntegrality} constructs a counterexample on the same
resolution $Y_+$ used here. Section~\ref{sec:integrality} explains
the distinction from the integer elementary expansion of
Ionel--Parker and records the scope of the transfer statement.
The specialization theorem and the compact calculation use only
the rational universal algebra.
